\documentclass[11pt,a4paper]{article}
\usepackage[utf8]{inputenc}
\usepackage[T1]{fontenc}
\usepackage{xcolor}
\usepackage[a4paper,margin=25mm]{geometry}
\usepackage{amsmath,amssymb,amsthm,mathtools,mathrsfs}
\usepackage{array,booktabs,tabularx,longtable}
\usepackage[expansion=false]{microtype}
\usepackage[colorlinks=true,linkcolor=blue,citecolor=blue,urlcolor=blue,hypertexnames=false]{hyperref}
\hypersetup{pdftitle={A zero-free half-plane for the Riemann zeta function},pdfsubject={An exposition of a cubic theta argument, matched to a Lean formalization}}
\allowdisplaybreaks[2]
\setlength{\emergencystretch}{3em}
\setcounter{tocdepth}{1}
\numberwithin{equation}{section}
\numberwithin{figure}{section}
\newtheorem{theorem}{Theorem}[section]
\newtheorem{proposition}[theorem]{Proposition}
\newtheorem{lemma}[theorem]{Lemma}
\newtheorem{corollary}[theorem]{Corollary}
\theoremstyle{definition}
\newtheorem{definition}[theorem]{Definition}
\newtheorem{slogan}[theorem]{Slogan}
\theoremstyle{remark}
\newtheorem{remark}[theorem]{Remark}
\newcommand{\Q}{\mathbb Q}
\newcommand{\R}{\mathbb R}
\newcommand{\NK}{\mathrm N_{K/\Q}}
\newcommand{\C}{\mathbb C}
\newcommand{\mf}{\mathfrak}
\newcommand{\OO}{\mathcal O}
\newcommand{\E}{\mathcal E}
\newcommand{\ind}{\mathbf1}
\newcommand{\sfstar}{\sideset{}{^*}\sum}
\newcommand{\rad}{\operatorname{rad}}
\newcommand{\Rea}{\operatorname{Re}}
\newcommand{\dd}{\,d}
\newcommand{\precbd}{\preccurlyeq}


\title{A zero-free half-plane for the Riemann zeta function:\\ $\zeta(s)\ne0$ for $\Rea s>23/24$}
\author{An exposition of OpenAI's \emph{The Quasi-Riemann Hypothesis},\\ matched to a Lean formalization}
\date{October 9, 2026}

